% ════════════════════════════════════════════
% 8. 结论
% ════════════════════════════════════════════
\section{结论与建议}

图\ref{fig:overview}给出全文结论的图形摘要：四个象限分别是付费带来的净效应 40.1\%、视频相对图片的倍数 4.90、模型的解释力 $R^2=0.207$、共线处理前后的最大方差膨胀因子 143.0 与 3.23，分别对应第 4 章与第 5 章的数值。

\subsection{结论汇总}

表\ref{tab:summary}按“分析任务—方法—核心结果”汇总全文，四项任务各自的判据与数值在正文对应章节可复核。

\begin{table}[htbp]
  \centering
  \caption{分析任务、方法与核心结果}
  \label{tab:summary}
  \zihao{-5}\setlength{\tabcolsep}{4pt}\renewcommand{\arraystretch}{1.05}
  \begin{tabular}{>{\centering\arraybackslash}p{0.16\textwidth}>{\centering\arraybackslash}p{0.30\textwidth}>{\centering\arraybackslash}p{0.44\textwidth}}
    \toprule
    分析任务 & 采用方法 & 核心结果 \\
    \midrule
    数据理解 & 结构核查、类型判定、偏度峰度与三项共线判据 & 帖级截面；$X_1$ 与月份 Spearman 0.996、辅助回归 $R^2=0.9917$、月内方差占比 0.83\%，三项判据同时成立，两变量不可同时识别 \\
    清洗与整理 & 众数填补、独热编码、$\log(1+Y_1)$ 变换、时段合并 & 缺失 1 条填补后无缺失；偏度由 8.33 降至 0.73、峰度由 93.05 降至 0.82；24 小时合并为凌晨、上午、下午晚间三组 \\
    回归预测 & OLS + HC3，五个候选模型，十折交叉验证 & M2 最优：$R^2=0.207$、调整 $R^2=0.165$、CV $R^2=0.101$、CV RMSE 1.101；最大 VIF 由 143.0 降至 3.23 \\
    效应解释 & Duan 烟雾还原、反事实模拟、Shapley 分解、六设定稳健性 & 付费提升 40.1\%（区间 10.8\%--77.2\%，代表性画像多 16,142 次）；视频 ×4.90（区间 2.23--10.79）；主题 2、3 分别低 38.0\% 与 27.0\%；类型贡献 37.1\%、月份 28.0\%、可解释 20.7\% \\
    \bottomrule
  \end{tabular}
\end{table}

三条可执行的投放结论是：付费推广在控制内容与时间后仍带来约四成曝光增量，且该结论在六种设定下方向一致；内容载体比投放更关键，视频帖的相对曝光是图片帖的近五倍，但样本量决定了这一倍数只能作为方向性依据；内容主题决定付费的边际收益，主题 1 的付费回报显著、主题 3 几乎没有，投放预算应优先配置给主题 1 的帖子。

\subsection{优点与适用边界}

\begin{itemize}
  \item 识别策略先于模型设定：共线判别放在建模之前，用三项可复核判据决定谁进入模型，避免了“先回归再看 VIF”的循环论证；剔除粉丝规模后最大 VIF 由 143.0 降至 3.23，所有系数进入可解释区间。
  \item 推断口径与残差诊断匹配：Breusch--Pagan 与 Jarque--Bera 均拒绝同方差与正态假设，因此全文使用 HC3 稳健标准误，并用 Cook 距离单独检查影响点，而不是在异方差存在时报告经典标准误。
  \item 适用范围：结论建立在单品牌、单期抓取、帖级截面数据上，因变量为 $\log(1+Y_1)$，系数解释为相对变化。跨品牌或跨平台推广时，粉丝规模的累积速度与算法分发规则不同，月份固定效应吸收的内容也不同，40.1\% 与 4.90 这两个数值需要重新标定；若需要绝对曝光量预测（而非相对比较），交叉验证 $R^2=0.101$ 的水平不足以支撑单帖预测。
  \item 样本边界：视频帖 7 条、链接帖 22 条，其系数区间宽度是图片帖的十倍以上；这三类结果的解释限定为“方向明确、幅度不确定”。
\end{itemize}

\subsection{改进方向}

\begin{itemize}
  \item 建模口径：以月份固定效应替代粉丝规模是本样本下的妥协。若取得同一主页的连续多期面板数据，可把粉丝规模与月份的共线拆开，用组内变换或动态面板估计粉丝规模的独立弹性。
  \item 模型扩展：曝光为计数型结果，可尝试负二项或 Gamma 回归直接建模原始尺度，与本文的对数线性设定互为稳健性对照；若引入帖子文本与图片特征，可用梯度提升树补充非线性与交互效应。
  \item 数据补充：当前缺少帖子发布时间戳与投放预算，无法区分“付费带来的增量”与“运营选择性投放”（高潜力帖子更可能被投放）。补充预算与投放时点后可做倾向得分匹配或双重差分，把选择效应从处理效应中剥离。
  \item 评估方式：交叉验证 $R^2=0.101$ 说明模型适合解释变量方向而非逐帖预测。若要用于投放决策，应改用分位数损失或提升度曲线（uplift）评估，并做 A/B 实验验证付费的真实增量。
\end{itemize}
